\documentclass[10pt,a4paper]{amsart}
\usepackage[margin=19mm]{geometry}
\usepackage{amsmath,amssymb,mathtools}
\usepackage[scr=rsfs,cal=euler]{mathalfa}
\usepackage{xfrac}
\usepackage[unicode,hidelinks,bookmarks=false]{hyperref}
\newcommand{\ie}{i.e.\ }
\newcommand{\cf}{cf.\ }
\newcommand{\resp}{respectively\ }
\newcommand{\Subsection}{\subsection}
\providecommand{\nobreakdash}{\nobreak\hbox{-}}
\setlength{\emergencystretch}{2em}
\multlinegap0pt
\usepackage{amssymb}
\usepackage{tikz}
\usepackage{float}
\usepackage[shortlabels]{enumitem}
\newtheorem{prop}{Proposition}
\def \Ppicc {{\scriptscriptstyle P}}
\newcommand{\R}{{\mathcal R}}
\def \Thetapicc {{\scriptscriptstyle \Theta}}
\def \lieg{\mathfrak{g}}
\def \liegRpP{\mathfrak{g}^{\raise1pt\hbox{$ \scriptscriptstyle \mathcal{R},p $}}_\Ppicc}
\def \liegRpPa{\mathfrak{g}^{\raise1pt\hbox{$ \scriptscriptstyle \mathcal{R},p $}}_{\Ppicc_{\raise-3pt\hbox{$ \scriptscriptstyle \!\! A $}}}}
\def \lieh{\mathfrak{h}}
\title{Multiparameter Quantum Supergroups, Deformations and Specializations}
\author{Gast\'on Andr\'es Garc\'ia, Fabio Gavarini, Margherita Paolini}
\date{Source: arXiv:2410.22549v2 (CC BY 4.0)}
\begin{document}
\maketitle

\noindent\emph{Source and licence.} Multiparameter Quantum Supergroups, Deformations and Specializations. Version arXiv:2410.22549v2 (CC BY 4.0), distributed by arXiv under the Creative Commons Attribution 4.0 International licence, http://creativecommons.org/licenses/by/4.0/, as recorded in the arXiv licence metadata for this exact version and shown on the abstract page https://arxiv.org/abs/2410.22549v2. Full licence notice: \texttt{licenses/arxiv-2410.22549-CC-BY-4.0.txt}.\par\smallskip

\noindent\textit{Editorial scope.} Counted unit: the article's prop liegRP-twist (source0.tex lines 2901--2928). The article's complete original proof of it is delivered with it (source0.tex lines 2930--2985, one \texttt{proof} environment of 56 lines). No mathematics is written, repaired or re-derived: the statement and the proof are the article's own lines, in the article's own notation, and no retained line is substituted or re-flowed.  Retained uncounted as the source-local context the proof needs: lines 2316--2320; lines 2869--2873.  Nothing was omitted from any retained span, so the package carries no omission record.  No retained line contains a citation, so the wrapper carries no bibliography and no bibliography entry.  No retained line contains a figure environment, an image-inclusion command or drawing code, so no image is delivered and the package declares no asset dependency.  Editorial formatting only, in addition: an AMS \texttt{amsart} preamble that restates the article's own declarations and macros used by the retained lines, namely the environments \texttt{prop} on separate counters, together with the standard packages those lines need.  The retained environments print in this document as Proposition 1 in order of appearance, whereas the article prints the same items with the numbering of its own full document.  The complete native source of the article as distributed in the arXiv e-print for this exact version is bundled unchanged as \texttt{sources/167638/source0.tex}.
\begin{equation}  \label{eq: def_twist-delta}
%
  \delta^{\,c}(x)  \; := \;  \delta(x) - (\partial c)(x)  \; = \;  \delta(x) - x\,.\,c   \qquad  \forall \;\; x \in \lieg
%
\end{equation}

\begin{equation}  \label{eq: def_twist_Lie-bialg}
  j_\Thetapicc  \; := \;  {\textstyle \sum_{g,k=1}^t} \theta_{gk} \, H_g \otimes H_k  \;
  \in \; \lieh \otimes \lieh  \; \subseteq \;  \lieg \otimes \lieg
%
\end{equation}

\begin{prop}  \label{prop: Yamane's twist-liegRPa = liegRP-twist}
%
   With assumptions as above, the following hold:
%
 \vskip5pt
%
   (a)\;  there exists a Lie superalgebra isomorphism
%
  $$  \phi^\Thetapicc_\Ppicc : \lieg_{P_\Theta}^{\R_\Thetapicc,p} \, {\buildrel \cong \over
 {\lhook\joinrel\relbar\joinrel\relbar\joinrel\relbar\joinrel\twoheadrightarrow}} \,
{\Big( \liegRpPa \Big)}^{j_\Thetapicc} \; ,  \qquad
%
  E_i \, \mapsto \, E_i \, , \; T \, \mapsto \, T \, , \; F_i \, \mapsto \, F_i \quad
  (\; i \in I \, , \, T \in \lieh \,)  $$
%
 \vskip1pt
%
   (b)\;  there exists a unique structure of Lie super\textsl{bi}algebra onto  $ \, \lieg_{P_\Theta}^{\R_\Thetapicc,p} \, $  such that the map  $ \phi^\Thetapicc_\Ppicc $
   in (a) above is an isomorphism of Lie super\textsl{bi}algebras;
%
 \vskip5pt
%
   (c)\;  the Lie cobracket of the Lie super\textsl{bi}algebra structure of  $ \, \lieg_{P_\Theta}^{\R_\Thetapicc,p} \, $  mentioned in (b) above is given by the formulas
   (for all  $ \, T \in \lieh \, $,  $ \, i \in I $)
%
  $$  \delta\big(\,T\big) \, = \, 0  \quad ,   \qquad  \delta\big(E_i\big) \, = \, 2 \; T^+_{\Thetapicc,i} \wedge E_i  \quad ,   \qquad  \delta\big(F_i\big) \, = \,  2 \; T^-_{\Thetapicc,i} \wedge F_i  $$
%
\end{prop}

\begin{proof}
%
 \textit{(a)}\,  By construction, the Lie algebra structure in  $ {\big( \liegRpPa \big)}^{j_\Thetapicc} $  is the same of  $ \liegRpPa \, $,
 and the latter only depends on the $ \alpha_j $'s  and the sums
 $ \, S_j := 2^{-1}\big(T_j^+ \! + T_j^-\big) \, $  ($ \, j \in I \, $).
 Now, both the  $ \alpha_j $'s  and the  $ S_j $'s  \textsl{do not change\/}
 (by construction)  when we pass from  $ \liegRpPa $  to  $ \lieg_{P_\Theta}^{\R_\Thetapicc,p} $
 or viceversa; therefore, the formulas in the claim   --- mapping each generator of  $
 \lieg_{P_\Theta}^{\R_\Thetapicc,p} $  onto the same name generator of
 $ \, \liegRpPa = {\big( \liegRpPa \big)}^{j_\Thetapicc} \, $  ---
 actually do provide an isomorphism  $ \phi^\Thetapicc_\Ppicc $  of Lie superalgebras,
 as claimed.
%
 \vskip5pt
%
   \textit{(b)}\,  By pull-back through  $ \phi^\Thetapicc_\Ppicc $  one gets onto
   $ \lieg_{P_\Theta}^{\R_\Thetapicc,p} $  the unique Lie cobracket which makes the former
   into a Lie superbialgebra such that the map  $ f^\Thetapicc_\Ppicc $
   is also an isomorphism of Lie super\textsl{bi\/}algebras.
%
 \vskip5pt
%
   \textit{(c)}\,  For the toral twist
%
  $ \; j_\Thetapicc := {\textstyle \sum_{g,k=1}^t} \theta_{gk} \, H_g \otimes H_k \; $
%
 in  \eqref{eq: def_twist_Lie-bialg},  formula  \eqref{eq: def_twist-delta}  yields
%
  $$  \delta^{j_\Thetapicc}(x)  \; := \;  \delta(x) - x.j_\Thetapicc  \; = \;
  \delta(x) - {\textstyle \sum_{g,k=1}^t} \theta_{gk} \,
  \big( \big[x,H_g\big] \otimes H_k + H_g \otimes \big[x,H_k\big] \big)  $$
%
 for all  $ \, x \in \liegRpPa \, $.  Now take  $ \, x := E_\ell \, $  ($ \, \ell \in I \, $):  \,then computations give
%
  $$  \displaylines{
   \delta^{j_\Thetapicc}(E_\ell)  \; := \;  \delta(E_\ell) - {\textstyle \sum_{g,k=1}^t} \, \theta_{gk} \,
   \big( \big[E_\ell \, , H_g\big] \otimes H_k + H_g \otimes \big[E_\ell \, , H_k\big] \big)  \; =   \hfill  \cr
   \quad   = \;  \delta(E_\ell) - {\textstyle \sum_{g,k=1}^t} \, \theta_{gk} \,
   \big(\! -\alpha_\ell(H_g) E_\ell \otimes H_k - H_g \otimes \alpha_\ell(H_k) E_\ell \,\big)  \; =   \hfill  \cr
   = \;  T^+_\ell \otimes E_\ell - E_\ell \otimes T^+_\ell + {\textstyle \sum_{g,k=1}^t} \, \theta_{gk} \,
   \big(\, \alpha_\ell(H_g) E_\ell \otimes H_k + \alpha_\ell(H_k) H_g \otimes E_\ell \,\big)  \; =   \quad  \cr
   \hfill   = \;  \Big(\, T^+_\ell + {\textstyle \sum_{g,k=1}^t} \, \theta_{kg} \, \alpha_\ell(H_g) H_k \Big) \otimes E_\ell \, -
   \, E_\ell \otimes \Big(\, T^+_\ell + {\textstyle \sum_{g,k=1}^t} \, \theta_{kg} \, \alpha_\ell(H_g) H_k \Big)  \; =  \cr
   \hfill   = \;  T^+_{\Thetapicc,\ell} \otimes E_\ell \, - \, E_\ell \otimes T^+_{\Thetapicc,\ell}  \; = \; 2 \;
   T^+_{\Thetapicc,\ell} \wedge \! E_\ell  }  $$
%
 hence in short we get  $ \; \delta^{j_\Thetapicc}(E_\ell) \, = \, 2 \; T^+_{\Thetapicc,\ell} \wedge \! E_\ell  \; $.
 Similar computations give
%
  $$
 \delta^{j_\Thetapicc}(T) = 0  \; ,   \!\quad
 \delta^{j_\Thetapicc}(F_\ell) = 2 \, T^-_{\Thetapicc,\ell} \wedge \! F_\ell  \; ,   \quad  \forall \; \ell \in I \, , \, T \!\in \lieh  $$
%
 Therefore, pulling back through  $ \phi_\Ppicc^\Thetapicc \, $  onto  $ \liegRpP $  the Lie cobracket of  $ {\Big( \liegRpPa \Big)}^{j_\Thetapicc} $,  \,we find that this Lie cobracket is described by the formulas in claim  \textit{(c)},  q.e.d.
%
\end{proof}

\end{document}
